% Part: normal-modal-logic
% Chapter: syntax-and-semantics
% Section: relational-models

\documentclass[../../../include/open-logic-section]{subfiles}

\begin{document}

\olfileid{nml}{syn}{rel}

\olsection{সম্পর্কমূলক মডেল}

স্বাভাবিক মোডাল যুক্তিবিদ্যার অর্থতত্ত্বের মূল ধারণা হলো \emph{সম্পর্কমূলক
মডেল}। এটি জগৎগুলির একটি সেট নিয়ে গঠিত; জগৎগুলি একটি দ্বিপদী
``অভিগম্যতা সম্পর্ক'' দ্বারা যুক্ত। আরও থাকে এমন একটি মাননির্ধারণ, যা
ঠিক করে কোন !!{propositional variable}s কোন জগতে ``সত্য'' ধরা হবে।

\begin{defn}
  মৌলিক মোডাল ভাষার একটি \emph{মডেল} হলো ত্রয়ী $\mModel{M}
  = \tuple{W, R, V}$, যেখানে
  \begin{enumerate}
  \item $W$ হলো ``জগৎ''-এর একটি অশূন্য সেট,
  \item $R$ হলো $W$-এর ওপর একটি দ্বিপদী অভিগম্যতা সম্পর্ক, এবং
  \item $V$ একটি অপেক্ষক, যা প্রতিটি !!{propositional
    variable}~$p$-কে সম্ভাব্য জগৎগুলির একটি সেট $V(p)$ নির্ধারণ করে।
  \end{enumerate}
  $Rww'$ সত্য হলে বলি $w$ থেকে $w'$ \emph{অভিগম্য}। $w \in V(p)$ হলে
  বলি $w$-তে $p$ \emph{সত্য}।
\end{defn}

সম্পর্কমূলক অর্থতত্ত্বের বড় সুবিধা হলো \olref{fig:simple}-এর মতো সহজ
চিত্রে মডেল দেখানো যায়। জগৎগুলি বিন্দু দিয়ে দেখানো হয়; $w$ থেকে
$w'$-এ তির থাকলেই এবং কেবল তখনই $w$ থেকে $w'$ অভিগম্য। উপরন্তু
$w \in V(p)$ হলে বিন্দুটিতে (জগৎটিতে)~$\mTrue{p}$ এবং অন্যথায়
\mFalse{p} লেখা হয়। \olref{fig:simple}-এ দেখানো মডেলে
$W = \{w_1, w_2, w_3\}$, $R = \{\tuple{w_1, w_2},\tuple{w_1,
  w_3}\}$, $V(p) = \{w_1, w_2\}$, এবং $V(q) = \{w_2\}$।

\begin{figure}
  \begin{center}
    \begin{tikzpicture}[modal]
      \node[world] (w1) [label={[align=right]right:\mTrue{p}\\ \mFalse{q}}]{$w_1$};
      \node[world] (w2) [label={[align=right]right:\mTrue{p}\\ \mTrue{q}},
        above right=of w1]{$w_2$};
      \node[world] (w3) [label={[align=right]right:\mFalse{p}\\ \mFalse{q}},
        below right=of w1] {$w_3$};
      \draw[->] (w1) to (w2);
      \draw[->] (w1) to (w3);
    \end{tikzpicture}
  \end{center}
  \caption{একটি সরল মডেল।}
  \ollabel{fig:simple}
\end{figure}

\end{document}
